\subsection{Структурно-функціональна модель підготовки майбутніх вчителів інформатики до розробки імерсивних освітніх ресурсів}\label{subsec:2.2}

У сучасному науково-педагогічному дискурсі професійне становлення фахівця осмислюється через побудову структурно впорядкованих моделей, що відображають логіку формування професійних умінь і компетентностей. Моделювання в педагогіці розглядається як загальнонауковий метод пізнання, що забезпечує опосередковане вивчення педагогічних об'єктів шляхом створення моделей і перенесення результатів дослідження з моделі на реальний освітній процес (С.~Гончаренко~\cite{Goncharenko2011Pedagogy}, О.~Дубасенюк~\cite{Dubaseniuk2008Modeling}, А.~Кузьмінський~\cite{Kuzminskyi2012Pedagogy}, Є.~Лодатко~\cite{Lodatko2011Modeling} тощо).

Сучасне осмислення моделювання зумовило увагу до класифікації педагогічних моделей, покладеної в основу системних досліджень, що уможливило виокремлення їх навчальних і науково-дослідницьких типів, де прогностичні варіанти відіграють провідну роль у підготовці фахівців до прогнозованих умов діяльності~\cite{Lodatko2011Modeling}. Вивчення наукового доробку вчених дозволило нам зазначити тотожність підходів, де вихідним положенням слугує поділ моделей за сферою застосування, забезпечуючи логічний перехід від констатуючих до імітаційних форм, здебільшого орієнтованих на експериментальне вивчення педагогічних систем. Подібна систематизація, на думку науковців, забезпечує глибше розуміння динаміки професійної освіти.

Вітчизняні педагогічні дослідження пропонують різні варіанти таких моделей, які відрізняються набором структурних компонентів і функціональних зв'язків між ними. Водночас систематичний огляд Q.~Wang та Y.~Li~\cite{Wang2024XRTeacherEd}, який охопив 52 публікації щодо підготовки вчителів засобами VR/AR/MR, засвідчив нагальну потребу в системному проєктуванні підготовки на основі теоретичних моделей (<<theory-driven training design>>). Фреймворк структурного IVR-навчання F.~Mulders та ін.~\cite{Mulders2020Framework} демонструє міжнародний підхід до побудови таких моделей із поєднанням когнітивних, конструктивістських і технологічних вимірів. Отже, розроблення структурно-функціональної моделі підготовки до створення ІОР відповідає актуальному глобальному запиту.

У процесі моделювання професійної підготовки більшість науковців виокремлюють структурні блоки/етапи, що забезпечують цілісність і логічну впорядкованість відповідних моделей. Так, І.~Дичківська~\cite{Dychkivska2012Innovative} у дослідженні інноваційних педагогічних технологій обґрунтовує модель, яка включає цільовий, змістовий, організаційно-технологічний та результативний блоки. За позицією П.~Мішри та М.~Коелера~\cite{Gurevych2015Integration}, концепція TPACK вибудовується як цілісна система взаємопов'язаних компонентів, де інтеграція технологічних, педагогічних та змістових знань визначає ефективність моделі підготовки вчителя; емпірично доведено, що TPACK-орієнтовані моделі підготовки вчителів до створення AR-ресурсів забезпечують суттєве зростання відповідних компетентностей~\cite{BeldaMedina2022TPACK}. С.~Гончаренко~\cite{Goncharenko2011Pedagogy} описує модель педагогічного дослідження, що складається з організаційного, змістового, процесуального та оцінно-результативного етапів. У процесі розроблення методології проєктування імерсивних освітніх ресурсів С.~Семеріков та співавтори~\cite{Semerikov2022IERDesign} виокремили взаємопов'язані компоненти: цільовий, теоретико-методологічний, організаційно-методичний та діагностико-результативний. Схожої логіки дотримуються С.~Литвинова, Н.~Сороко та С.~Баценко~\cite{Lytvynova2023Monograph}, які визначають мотиваційно-цільовий, змістовий, організаційно-процесуальний, контрольно-оціночний та результативний блоки. Є.~Лодатко пропонує виокремлювати мотиваційно-діагностичний, навчально-професійний та рефлексивно-оцінний модулі~\cite{Lodatko2011Modeling}. Водночас динамічна модель ADDIE, адаптована для VR-навчання, визначає <<інтерактивний дизайн зворотного зв'язку>> як найвпливовіший фактор ефективності імерсивних навчальних систем~\cite{Li2024ADDIE}.

Опрацювання сучасної науково-педагогічної літератури дало підстави стверджувати, що конструювання структурно-функціональної моделі професійної підготовки має ґрунтуватися на внутрішній узгодженості мети, завдань, принципів, підходів, педагогічних умов та критеріїв її реалізації.

У процесі розроблення моделі підготовки майбутніх учителів інформатики до розробки імерсивних освітніх ресурсів ми спиралися на наявні в сучасному науково-педагогічному дискурсі теоретично обґрунтовані моделі професійної підготовки вчителів інформатики за різними напрямами. Зокрема, у працях П.~Мішри та М.~Коелера обґрунтовано концепцію TPACK, яка акцентує взаємозв'язок технологічного, змістового та педагогічного компонентів професійної підготовки вчителя~\cite{Gurevych2015Integration}. У зазначених дослідженнях підкреслюється значення системності, поетапності та методичної доцільності організації освітнього процесу, а також орієнтація на формування професійних умінь, релевантних до реальних умов педагогічної діяльності.

Адаптуючи положення окреслених наукових підходів з урахуванням специфіки імерсивних освітніх ресурсів, їх технологічної складності та педагогічної функціональності, нами запропоновано власну модель підготовки майбутніх учителів інформатики, зорієнтовану на поєднання технологічної, дидактичної та проєктної складових фахової підготовки.

Спираючись на окреслені теоретико-методологічні положення (розділ~2.1.), нами розроблено структурно-функціональну модель підготовки майбутніх учителів інформатики до розробки ІОР у процесі фахової підготовки. Така модель розглядається як цілісна система, у межах якої окремі структурні компоненти та функціональні зв'язки між ними забезпечують цілеспрямованість і логіку процесу формування професійних умінь. У ході її конструювання враховано сукупність ключових структурних компонентів, які в сучасних педагогічних дослідженнях розглядаються як методологічно значущі для формування професійних умінь фахівців означеного профілю та забезпечення їх готовності до педагогічно вмотивованого використання ІТ у майбутній професійній діяльності. Усі складники моделі є взаємопов'язаними, взаємозумовленими й займають визначене місце в загальній структурі фахової підготовки. Виходячи з цього, структурно-функціональна модель ІОР освітніх ресурсів передбачає поетапну організацію освітнього процесу, що охоплює \emph{мотиваційно-цільовий,} \emph{теоретико-методологічний, змістовий, організаційно-технологічний та рефлексивно-оцінний блоки.}

Кожен структурний блок моделі відповідає конкретним медіаторним факторам когнітивно-афективної моделі імерсивного навчання (CAMIL)~\cite{Makransky2021CAMIL} та інтегрує визначені компоненти TPACK~\cite{Gurevych2015Integration}: мотиваційно-цільовий блок активізує фактори мотивації та саморегуляції (компоненти PK, CK); теоретико-методологічний --- фактор когнітивного навантаження (TK+PK+CK); змістовий --- фактори втілення й агентності (CK, TK); організаційно-технологічний --- фактори агентності й присутності (TK, TPK); рефлексивно-оцінний --- фактор саморегуляції (TPK, TPACK). Послідовність блоків відтворює трифазну педагогічну структуру (теоретична → практична → апробаційна), закладену в 6-компонентній моделі проєктування ІОР~\cite{Semerikov2022IERDesign}. Як засвідчує систематичний огляд~\cite{Mallek2024SynthesizingVR}, узгодження педагогічного дизайну з когнітивними теоріями є ключовою передумовою ефективності імерсивного навчання. Зведення блоків моделі, факторів CAMIL, компонентів TPACK та етапів 6-компонентної моделі подано на рис.~\ref{fig:model_alignment}.

\begin{figure}[!h]
\centering
%\noindent\resizebox{\textwidth}{!}{%
%\begin{tikzpicture}
%\tikzset{
%  hdr/.style={fill=black!15, font=\footnotesize\bfseries, minimum height=0.7cm, text centered, inner sep=3pt},
%  cel/.style={font=\footnotesize, minimum height=1.0cm, text centered, inner sep=3pt, text width=#1, align=center},
%}
%
%% Column widths
%\def\colA{4.6cm}
%\def\colB{3.8cm}
%\def\colC{3.2cm}
%\def\colD{3.6cm}
%
%% Header row
%\node[hdr, text width=\colA] (ha) at (0,0) {Блок моделі};
%\node[hdr, text width=\colB, right=-\pgflinewidth of ha] (hb) {Фактори CAMIL};
%\node[hdr, text width=\colC, right=-\pgflinewidth of hb] (hc) {Компоненти TPACK};
%\node[hdr, text width=\colD, right=-\pgflinewidth of hc] (hd) {Етап 6-комп. моделі};
%
%% Row 1: Мотиваційно-цільовий
%\node[cel=\colA, below=-\pgflinewidth of ha] (a1) {Мотиваційно-цільовий};
%\node[cel=\colB, right=-\pgflinewidth of a1] (b1) {мотивація,\\саморегуляція};
%\node[cel=\colC, right=-\pgflinewidth of b1] (c1) {PK, CK};
%\node[cel=\colD, right=-\pgflinewidth of c1] (d1) {Аналіз};
%
%% Row 2: Теоретико-методологічний
%\node[cel=\colA, below=-\pgflinewidth of a1] (a2) {Теоретико-\\методологічний};
%\node[cel=\colB, right=-\pgflinewidth of a2] (b2) {когнітивне\\навантаження};
%\node[cel=\colC, right=-\pgflinewidth of b2] (c2) {TK+PK+CK};
%\node[cel=\colD, right=-\pgflinewidth of c2] (d2) {Вимоги};
%
%% Row 3: Змістовий
%\node[cel=\colA, below=-\pgflinewidth of a2] (a3) {Змістовий};
%\node[cel=\colB, right=-\pgflinewidth of a3] (b3) {втілення,\\агентність};
%\node[cel=\colC, right=-\pgflinewidth of b3] (c3) {CK, TK};
%\node[cel=\colD, right=-\pgflinewidth of c3] (d3) {Підбір ПЗ +\\планування};
%
%% Row 4: Організаційно-технологічний
%\node[cel=\colA, below=-\pgflinewidth of a3] (a4) {Організаційно-\\технологічний};
%\node[cel=\colB, right=-\pgflinewidth of a4] (b4) {агентність,\\присутність};
%\node[cel=\colC, right=-\pgflinewidth of b4] (c4) {TK, TPK};
%\node[cel=\colD, right=-\pgflinewidth of c4] (d4) {Розробка +\\впровадження};
%
%% Row 5: Рефлексивно-оцінний
%\node[cel=\colA, below=-\pgflinewidth of a4] (a5) {Рефлексивно-оцінний};
%\node[cel=\colB, right=-\pgflinewidth of a5] (b5) {саморегуляція};
%\node[cel=\colC, right=-\pgflinewidth of b5] (c5) {TPK, TPACK};
%\node[cel=\colD, right=-\pgflinewidth of c5] (d5) {Оцінювання};
%
%% Draw grid borders
%\foreach \r in {a,b,c,d} {
%  \foreach \i in {1,2,3,4,5} {
%    \draw (\r\i.north west) rectangle (\r\i.south east);
%  }
%}
%% Header borders
%\draw (ha.north west) rectangle (ha.south east);
%\draw (hb.north west) rectangle (hb.south east);
%\draw (hc.north west) rectangle (hc.south east);
%\draw (hd.north west) rectangle (hd.south east);
%
%% Outer frame
%\draw[thick] (ha.north west) rectangle (d5.south east);
%% Thick line below header
%\draw[thick] (ha.south west) -- (hd.south east);
%
%\end{tikzpicture}%
%}
{\small
	\begin{tabularx}{\linewidth}{|p{4.3cm} |l |c|X|}
		\hline
		\rowcolor{gray!20}
		\hfil \textbf{Блок моделі} & 
		\hfill \textbf{Фактори CAMIL} & 
		\textbf{\makecell{Компоненти \\TPACK}} & 
		\hfil \textbf{Етап 6-комп. моделі} \\
		\hline
		
		Мотиваційно-цільовий 
		& мотивація, саморегуляція 
		& PK, CK 
		& Аналіз \\
		\hline
		Теоретико-методологічний 
		& когнітивне навантаження 
		& TK+PK+CK 
		& Вимоги \\
		\hline
		Змістовий 
		& втілення, agentність 
		& CK, TK 
		& Підбір ПЗ + планування \\
		\hline
		Організаційно-технологічний 
		& агентність, присутність 
		& TK, TPK 
		& Розробка + впровадження \\
		\hline
		Рефлексивно-оцінний 
		& саморегуляція 
		& TPK, TPACK 
		& Оцінювання \\
		\hline
	\end{tabularx}
}
\caption{Відповідність блоків структурно-функціональної моделі факторам CAMIL, компонентам TPACK та етапам 6-компонентної моделі проєктування ІОР}
\label{fig:model_alignment}
\end{figure}

Зазначена структурно-функціональна модель підготовки майбутніх учителів інформатики до розробки імерсивних освітніх ресурсів, що відображає логіку взаємодії її основних блоків, подана на рис.~\ref{fig:2.3}.

\begin{figure}[!h]
\centering
\noindent\resizebox{\textwidth}{!}{%
\begin{tikzpicture}[node distance=0.4cm, every node/.style={font=\footnotesize}]

% ============================================================
% BLOCK 1: МОТИВАЦІЙНО-ЦІЛЬОВИЙ БЛОК
% ============================================================
\node[dia/header=14.5cm, minimum height=0.8cm] (h1) {МОТИВАЦІЙНО-ЦІЛЬОВИЙ БЛОК};

\node[dia/lrbox=14.5cm, below=0.1cm of h1, inner sep=3pt] (b1) {%
  \begin{minipage}[t]{5.0cm}\setlength{\parskip}{0pt}\setlength{\baselineskip}{10pt}%
  \textbf{\textit{Мета:}} формування готовності майбутніх учителів інформатики до розробки і педагогічно доцільного використання імерсивних освітніх ресурсів у професійній діяльності.
  \end{minipage}%
  \hfill
  \begin{minipage}[t]{8.7cm}\setlength{\parskip}{0pt}\setlength{\baselineskip}{10pt}%
  \textbf{\textit{Завдання:}}\\[1pt]
  \checkmark~розвиток стійкої мотивації до опанування ІТ як складника сучасної педагогічної практики;\\[0pt]
  \checkmark~забезпечення досягнення очікуваних результатів професійної підготовки, пов'язаних із проєктуванням і реалізацією ІОР;\\[0pt]
  \checkmark~формування професійної рефлексії майбутніх учителів інформатики щодо власної діяльності у сфері створення та застосування ІОР
  \end{minipage}%
};

\node[draw, thick, inner sep=0pt, fit=(h1)(b1)] (block1) {};

% ============================================================
% BLOCK 2: ТЕОРЕТИКО-МЕТОДОЛОГІЧНИЙ БЛОК
% ============================================================
\node[dia/header=14.5cm, minimum height=0.8cm, below=0.7cm of block1.south] (h2) {ТЕОРЕТИКО-МЕТОДОЛОГІЧНИЙ БЛОК};

\node[dia/lrbox=14.5cm, below=0.1cm of h2, inner sep=3pt] (b2) {%
  \begin{minipage}[t]{7.5cm}\setlength{\parskip}{0pt}\setlength{\baselineskip}{10pt}%
  \textbf{\textit{Дидактичні\,підходи:}} аксіологічний, діяльнісний, інноваційний, інтегративний, когнітивний, компетентнісний, конструктивістський, контекстний, рефлексивний, синергетичний і системний.
  \end{minipage}%
  \hfill
  \begin{minipage}[t]{6.2cm}\setlength{\parskip}{0pt}\setlength{\baselineskip}{10pt}%
  \textbf{\textit{Принципи\,навчання:}} науковість, інтерактивність, варіативність, практична спрямованість, інтегрованість, рефлексивність, креативність, безперервність.
  \end{minipage}%
};

\node[draw, thick, inner sep=0pt, fit=(h2)(b2)] (block2) {};
\draw[dia/darrow] (block1.south) -- (block2.north);

% ============================================================
% BLOCK 3: ЗМІСТОВИЙ БЛОК
% ============================================================
\node[dia/header=14.5cm, minimum height=0.8cm, below=0.7cm of block2.south] (h3) {ЗМІСТОВИЙ БЛОК};

\node[dia/lrbox=14.5cm, below=0.1cm of h3, inner sep=3pt] (b3) {%
  \begin{minipage}[t]{7.2cm}\setlength{\parskip}{0pt}\setlength{\baselineskip}{10pt}%
  \textbullet~навчальні дисципліни професійної підготовки (обов'язкові компоненти ОПП);\\[0pt]
  \textbullet~вибіркові дисципліни, спрямовані на поглиблене вивчення програмування, комп'ютерної графіки, моделювання, інтерактивних технологій, основ VR/AR/MR та дидактичного дизайну цифрових освітніх ресурсів.
  \end{minipage}%
  \hfill
  \begin{minipage}[t]{6.5cm}\setlength{\parskip}{0pt}\setlength{\baselineskip}{10pt}%
  \textbullet~навчально-дослідницька та проєктна діяльність здобувачів;\\[0pt]
  \textbullet~практична підготовка охоплює педагогічну та виробничу практики, у межах яких майбутні вчителі інформатики апробують розроблені ІОР у реальному або наближеному до реального освітньому середовищі.
  \end{minipage}%
};

\node[draw, thick, inner sep=0pt, fit=(h3)(b3)] (block3) {};
\draw[dia/darrow] (block2.south) -- (block3.north);

% ============================================================
% BLOCK 4: ОРГАНІЗАЦІЙНО-ТЕХНОЛОГІЧНИЙ БЛОК
% ============================================================
\node[dia/header=14.5cm, minimum height=0.8cm, below=0.7cm of block3.south] (h4) {ОРГАНІЗАЦІЙНО-ТЕХНОЛОГІЧНИЙ БЛОК};

\node[dia/lrbox=14.5cm, below=0.1cm of h4, inner sep=3pt] (b4) {%
  \begin{minipage}[t]{4.2cm}\setlength{\parskip}{0pt}\setlength{\baselineskip}{10pt}%
  \textbf{\textit{Форми:}} лекції, лабораторні та практичні заняття, самостійна робота здобувачів, тренінги, у тому числі комп'ютерні, проєктна діяльність, ділові ігри.
  \end{minipage}%
  \hfill
  \begin{minipage}[t]{5.6cm}\setlength{\parskip}{0pt}\setlength{\baselineskip}{10pt}%
  \textbf{\textit{Методи:}} пояснювально-ілюстративні, частково-пошукові, дослідницькі, евристичні, бесіди, дискусії, метод проєктів, інтерактивні/інноваційні, метод проблемного викладу, аналізу ситуацій.
  \end{minipage}%
  \hfill
  \begin{minipage}[t]{3.4cm}\setlength{\parskip}{0pt}\setlength{\baselineskip}{10pt}%
  \textbf{\textit{Засоби:}} словесні, наочні, технічні, засоби навчання, прикладне програмне забезпечення спеціального призначення.
  \end{minipage}%
};

\node[draw, thick, inner sep=0pt, fit=(h4)(b4)] (block4) {};
\draw[dia/darrow] (block3.south) -- (block4.north);

% ============================================================
% BLOCK 5: РЕФЛЕКСИВНО-ОЦІННИЙ БЛОК
% ============================================================
\node[dia/header=14.5cm, minimum height=0.8cm, below=0.7cm of block4.south] (h5) {РЕФЛЕКСИВНО-ОЦІННИЙ БЛОК};

\node[dia/lrbox=14.5cm, below=0.1cm of h5, inner sep=3pt] (b5) {%
  \begin{minipage}[t]{4.2cm}\setlength{\parskip}{0pt}\setlength{\baselineskip}{10pt}%
  \textbf{\textit{Компоненти:}} когнітивний; мотиваційно-ціннісний; операційно-діяльнісний; рефлексивний
  \end{minipage}%
  \hfill
  \begin{minipage}[t]{5.8cm}\setlength{\parskip}{0pt}\setlength{\baselineskip}{10pt}%
  \textbf{\textit{Критерії:}} мотиваційно-ціннісний; діяльнісно-практичний; методично-дидактичний; когнітивно-результативний
  \end{minipage}%
  \hfill
  \begin{minipage}[t]{3.2cm}\setlength{\parskip}{0pt}\setlength{\baselineskip}{10pt}%
  \textbf{\textit{Рівні:}} високий; середній; низький
  \end{minipage}%
};

\node[draw, thick, inner sep=0pt, fit=(h5)(b5)] (block5) {};
\draw[dia/darrow] (block4.south) -- (block5.north);

% ============================================================
% LEFT SIDEBAR: ПЕДАГОГІЧНІ УМОВИ (with arrows to each block)
% ============================================================
% Sidebar rectangle spanning all 5 blocks
% Use fit to auto-size, then overlay text
\coordinate (sbar mid) at ($(block1.north west)!0.5!(block5.south west) + (-1.8cm,0)$);
\node[draw, thick, rotate=90, anchor=center,
      minimum height=2.4cm,
      font=\footnotesize\bfseries, text centered, inner sep=2pt,
      minimum width=20cm]
  at (sbar mid) (sidebar) {ПЕДАГОГІЧНІ УМОВИ};

% Arrows from sidebar right edge to each block
\coordinate (sbar edge) at ($(sbar mid) + (1.2cm,0)$);
\draw[dia/arrow] (sbar edge |- block1) -- (block1.west);
\draw[dia/arrow] (sbar edge |- block2) -- (block2.west);
\draw[dia/arrow] (sbar edge |- block3) -- (block3.west);
\draw[dia/arrow] (sbar edge |- block4) -- (block4.west);
\draw[dia/arrow] (sbar edge |- block5) -- (block5.west);

% ============================================================
% BOTTOM: Результат
% ============================================================
\node[dia/lrbox=14.5cm, below=0.7cm of block5.south, inner sep=3pt] (result) {%
  \begin{minipage}[t]{14.1cm}\setlength{\parskip}{0pt}\setlength{\baselineskip}{10pt}%
  \textbf{\textit{Результат:}} готовність майбутніх учителів інформатики до реалізації методики розробки імерсивних освітніх ресурсів, що забезпечує інтеграцію технологічних, дидактичних і проєктних умінь у процесі фахової підготовки та їх застосування у реальних педагогічних ситуаціях.%
  \end{minipage}%
};

\draw[dia/arrow] (block5.south) -- (result.north);

\end{tikzpicture}%
}

\caption{Структурно-функціональна модель навчання розробки імерсивних освітніх ресурсів майбутніх учителів інформатики}
\label{fig:2.3}
\end{figure}

Подальший виклад спрямовано на розкриття змістового наповнення виокремлених структурних блоків моделі з урахуванням їхньої цільової спрямованості, функціонального призначення та методологічних орієнтирів.

\emph{\textbf{Мотиваційно-цільовий блок}} у структурно-функціональній моделі підготовки майбутніх учителів інформатики до розробки ІОР виконує вихідну, смислотвірну та спрямувальну функції. Блок передбачає обґрунтування соціального запиту на підготовку професійно компетентних майбутніх учителів інформатики, здатних до розробки та педагогічно доцільного використання ІОР, врахування потреб особистості у професійному саморозвитку, цифровій мобільності та адаптації до умов трансформації освітнього середовища. Реалізація зазначених положень здійснюється з урахуванням особливостей, принципів і завдань навчання, що зумовлюються метою моделі - формування готовності майбутніх учителів інформатики до розробки імерсивних освітніх ресурсів у межах фахової підготовки. Як засвідчує дослідження T.~Lee та ін.~\cite{Lee2022VRMaking}, досвід створення VR-ресурсів (VR-Making) сприяє формуванню технологічної готовності, навичок <<4C>> (критичне мислення, комунікація, колаборація, креативність) та усвідомлення педагогічних переваг імерсивних технологій. Водночас за результатами структурного моделювання (PLS-SEM, \textit{N}=278) самоефективність є найсильнішим предиктором готовності майбутніх учителів до впровадження VR-систем навчання~\cite{Xie2024UTAUT2VR}, що підтверджує значущість мотиваційного компонента в структурі моделі (фактор мотивації CAMIL).

Нормативно-правове підґрунтя розробленої моделі формується на основі положень Закону України «Про вищу освіту», проєкту стандарту вищої освіти за спеціальністю 014.09 «Середня освіта (Інформатика)» для першого (бакалаврського) рівня, а також чинних ОПП і робочих навчальних планів підготовки майбутніх учителів інформатики, що впроваджуються в закладах вищої освіти у межах відповідних освітніх програм бакалаврського рівня.

Саме в межах цього блоку закладаються ціннісно-мотиваційні основи професійної діяльності, формується усвідомлене ставлення здобувачів освіти до ІТ як до важливого ресурсу сучасної педагогічної практики, а також визначається загальне спрямування подальших етапів підготовки.

\emph{Метою} мотиваційно-цільового блоку є формування готовності майбутніх учителів інформатики до розробки і педагогічно доцільного використання ІОР у професійній діяльності. Мета постає не лише як очікуваний результат навчання, а як інтегративний орієнтир, що поєднує професійні, технологічні та педагогічні аспекти підготовки, узгоджуючи їх із сучасними вимогами до цифрової трансформації освіти.

До основних \emph{завдань} мотиваційно-цільового блоку належить розвиток стійкої мотивації майбутніх учителів інформатики до опанування ІТ як складника сучасної педагогічної практики; забезпечення орієнтації здобувачів освіти на досягнення очікуваних результатів професійної підготовки, пов'язаних із проєктуванням та реалізацією ІОР; формування професійної рефлексії щодо власної діяльності у сфері створення та застосування ІОР. Реалізація цих завдань сприяє усвідомленню значущості ІТ для підвищення якості навчання та розвитку інноваційного педагогічного мислення.

\emph{\textbf{Теоретико-методологічний блок}} у структурно-функціональній моделі підготовки майбутніх учителів інформатики до розробки ІОР виконує концептуально-обґрунтовувальну функцію. Саме в його межах визначаються наукові підходи та принципи, які задають методологічну рамку всього процесу професійної підготовки майбутніх учителів інформатики й забезпечують його внутрішню логічну узгодженість.

\emph{Метою} теоретико-методологічного блоку є формування цілісного науково-педагогічного підґрунтя підготовки майбутніх учителів інформатики, що забезпечує осмислене поєднання теоретичних знань про ІОР з методикою їх педагогічного проєктування та використання у професійній діяльності. У цьому контексті теоретико-методологічні засади виступають джерелом знань та орієнтиром для вибору адекватних способів організації навчання, форм і методів освітньої взаємодії.

До основних \emph{завдань} теоретико-методологічного блоку належить обґрунтування сукупності наукових підходів, що забезпечують ефективність підготовки майбутніх учителів інформатики до розробки ІОР; визначення принципів навчання, які відображають специфіку ІТ та педагогічної діяльності в умовах цифрової трансформації освіти; інтеграція положень педагогіки, психології, інформатики та цифрових технологій у єдину методологічну систему. Реалізація цих завдань сприяє формуванню методологічної культури майбутніх учителів та здатності свідомо застосовувати теоретичні знання у практичній діяльності.

Теоретико-методологічний блок спирається на сукупність дидактичних підходів (докладно у підрозділі~\ref{subsec:2.1}), зокрема аксіологічний, діяльнісний, інноваційний, інтегративний, когнітивний, компетентнісний, конструктивістський, контекстний, рефлексивний, синергетичний і системний, які забезпечують цілісне осмислення процесу підготовки майбутніх учителів інформатики до розробки ІОР та орієнтують його на активну, усвідомлену й професійно значущу діяльність. Застосування зазначених підходів дозволяє поєднати ціннісні орієнтири професійної підготовки з практичним досвідом проєктування імерсивних освітніх рішень, а також врахувати когнітивні, діяльнісні й рефлексивні аспекти становлення майбутнього педагога.

Систему принципів навчання у межах теоретико-методологічного блоку становлять принципи науковості, інтерактивності, варіативності, практичної спрямованості, інтеграції, рефлексивності, креативності та безперервності. Їх реалізація забезпечує методологічну обґрунтованість і педагогічну доцільність процесу підготовки, сприяє узгодженню теоретичних положень із практичними завданнями розробки ІОР та створює підґрунтя для послідовного переходу до змістового блоку моделі, зберігаючи наступність між цільовими орієнтирами й практичною підготовкою майбутніх учителів інформатики.

\emph{\textbf{Змістовий блок}} у структурно-функціональній моделі підготовки майбутніх учителів інформатики до розробки ІОР, спрямований на формування цілісного комплексу знань і уявлень, необхідних для усвідомленого проєктування та педагогічно доцільного використання технологій занурення. У цьому блоці конкретизовано навчальний зміст підготовки з урахуванням як технологічних, так і дидактичних вимог до ІОР.

\emph{Метою змістового блоку} є наповнення процесу підготовки структурованими та науково обґрунтованими знаннями про ІОР, принципи їх функціонування, можливості й обмеження використання в освітньому середовищі, а також про засади дидактичного дизайну імерсивних ресурсів. У цьому контексті визначається зміст навчальних дисциплін професійної підготовки, що входять до обов'язкових компонентів ОПП, а також вибіркових дисциплін, орієнтованих на поглиблене вивчення програмування, комп'ютерної графіки, моделювання, інтерактивних технологій, основ VR/AR/MR та дидактичного дизайну цифрових освітніх ресурсів. Така структура змісту забезпечує системне опанування теоретичних основ і технологічних інструментів, необхідних для створення імерсивних освітніх продуктів.

До \emph{основних завдань} змістового блоку належить інтеграція знань з інформатики, педагогіки, психології навчання та освітнього дизайну, формування у майбутніх учителів інформатики цілісного уявлення про імерсивний освітній ресурс як дидактично виважений продукт педагогічної діяльності, а не лише як технічну розробку.

За основу було використано ОПП Інформатика. Програмування КДПУ, де було виявлені ключові дисципліни циклу психолого-педагогічної, методичної та науково-предметної підготовки, зокрема «Методика навчання інформатики» (ОК 15), «Цифрові технології в освіті» (ОК 18), «Розробка цифрових освітніх ресурсів» (ОК 21), «Комп'ютерна графіка» (ОК 24), «Людино-машинна взаємодія» (ОК 32), а також навчальні й виробничі практики (ОК 36, ОК 37, ОК 40, ОК 42), формують комплекс знань, умінь і навичок, опанування яких забезпечує готовність майбутніх учителів інформатики до розробки імерсивних освітніх ресурсів.

Зазначені освітні компоненти безпосередньо спрямовані на розвиток інформаційно-цифрової компетентності (СК-03), проєктувальної компетентності (СК-09), інноваційної компетентності (СК-13) та здатності проєктувати й розробляти програмне забезпечення з використанням сучасних парадигм програмування (СК-16), що корелює з програмними результатами навчання ПР 02, ПР 03, ПР 09, ПР 13 та ПР 16. У центрі змістового блоку запропонованої структурно-функціональної моделі перебуває дидактичне наповнення й методичне забезпечення спеціально спроєктованого курсу \emph{\textbf{«Синтетична реальність в освіті»}}, орієнтованого на інтеграцію педагогічних, технологічних і дизайнерських компонентів та формування здатності до педагогічно вмотивованого проєктування, створення й використання імерсивних технологій у реальному освітньому середовищі. Структура курсу реалізує послідовність <<теоретична → практична → апробаційна>>, закладену в 6-компонентній моделі проєктування ІОР~\cite{Semerikov2022IERDesign}, де відбір змісту спирається на 23 ключові фактори динамічної моделі ADDIE для VR-навчальних систем~\cite{Li2024ADDIE}. Такий підхід <<навчання через дизайн>> (learning by design) підтверджено квазіекспериментально: X.-F.~Lin та ін.~\cite{Lin2024ARLearningByDesign} встановили, що AR-орієнтована проєктна діяльність суттєво підвищує мислення вищого порядку студентів педагогічних спеціальностей (N=128). Водночас дизайн навчального досвіду для AR потребує інтеграції принципів UX-дизайну з педагогічними цілями~\cite{Czerkawski2021LXD}, що враховано під час формування змісту курсу.

Вагоме місце у структурі змістового блоку займає навчально-дослідницька та проєктна діяльність здобувачів освіти, реалізація якої забезпечується в межах фахових і методичних дисциплін, насамперед курсу з методики навчання інформатики, розробки цифрових освітніх ресурсів, комп'ютерної графіки, людино-машинної взаємодії та програмування. У процесі виконання курсових і проєктних завдань студенти застосовують набуті знання й уміння для проєктування та створення власних імерсивних освітніх ресурсів, поєднуючи технологічні рішення з дидактичними цілями навчання. Особливого значення набуває практична підготовка, представлена навчальними та виробничими практиками, зокрема педагогічною та практикою в закладах освіти, під час яких майбутні вчителі інформатики здійснюють апробацію створених імерсивних ресурсів у реальному або наближеному до реального освітньому середовищі, оцінюють їх дидактичну доцільність, функціональність і відповідність освітнім результатам.

Змістовий блок зорієнтований на цілісне поєднання навчального матеріалу з майбутньою професійною діяльністю вчителя інформатики та враховує специфіку розробки імерсивних освітніх ресурсів. Його реалізація спирається на систему дидактичних підходів і принципів, що визначають логіку відбору, структурування й оновлення навчального змісту.

Реалізація змістового блоку створює методичне й дидактичне підґрунтя для переходу до \emph{\textbf{організаційно-технологічного}} блоку моделі, у межах якого засвоєні знання набувають прикладного характеру та трансформуються у практичні вміння і навички розробки, апробації та педагогічного використання імерсивних освітніх ресурсів у реальному або наближеному до реального освітньому середовищі.

\emph{Метою} організаційно-технологічного блоку є забезпечення практико-орієнтованої реалізації процесу підготовки майбутніх учителів інформатики до розробки імерсивних освітніх ресурсів, узгодження цілей і завдань фахової підготовки з вимогами освітніх стандартів, специфікою професійної діяльності вчителя інформатики та технологічними особливостями імерсивних освітніх рішень. У межах цього блоку конкретизуються завдання формування відповідних умінь, визначається місце імерсивних технологій у структурі освітньої програми, а також забезпечується взаємозв'язок між теоретичною підготовкою, практичною діяльністю та проєктною роботою студентів.

Організаційно-технологічний блок охоплює процесуально-діяльнісний складник, спрямовану на безпосереднє формування практичних умінь і досвіду проєктної діяльності майбутніх учителів інформатики. В межах цього блоку відбувається перехід від засвоєння теоретичних положень до їх реалізації у практиці створення імерсивних освітніх продуктів, що відповідають дидактичним цілям, навчальному контексту та потребам здобувачів освіти. Як засвідчує дослідження M.~Oberdörfer та ін.~\cite{Oberdorfer2021Interdisciplinary}, міждисциплінарна підготовка студентів педагогічних і технічних спеціальностей до проєктування VR/AR-навчальних середовищ забезпечує <<взаємну вигоду>> --- одночасний розвиток педагогічних і технологічних компетентностей. Організаційною основою практичної діяльності слугує цикл емпіричного навчання Д.~Колба~\cite{Kee2023Experiential}, де конкретний досвід (CE) → рефлексивне спостереження (RO) → абстрактна концептуалізація (AC) → активне експериментування (AE) реалізується як: лабораторна робота → рефлексія → теоретичне узагальнення → проєктне створення. Практичну реалізацію цих організаційних принципів забезпечує авторський курс <<Синтетична реальність в освіті>> (18 тижневих модулів), побудований на прогресивному ускладненні від базового WebAR (відстеження зображень MindAR) до інтегрованого ML-WebAR (TensorFlow.js)~\cite{Semerikov2024WebARML_CoSinE}, із використанням навчального посібника з WebAR~\cite{Shepiliev2020WebAR} як основного навчального ресурсу.

До основних завдань організаційно-технологічного блоку належить організація практико-орієнтованої діяльності студентів, залучення їх до виконання навчально-проєктних завдань, що передбачають застосування інструментів програмування, моделювання, інтерактивної візуалізації, прототипування та тестування імерсивних освітніх ресурсів. Важливого значення набуває розвиток умінь планувати й реалізовувати проєкти, працювати в команді, здійснювати рефлексію результатів власної діяльності, а також створення організаційних і технологічних умов для інтеграції аудиторної роботи з самостійною, командною, проєктною та дослідницькою діяльністю здобувачів освіти.

Організаційно-технологічний блок спрямований на активізацію навчальної діяльності здобувачів освіти та створення ними власного освітнього продукту в процесі фахової підготовки. Його реалізація забезпечує логічний перехід до діагностично-результативного блоку, орієнтованого на оцінювання й узагальнення досягнутих результатів підготовки.

\emph{\textbf{Рефлексивно-оцінний блок}} у структурно-функціональній моделі підготовки майбутніх учителів інформатики до розробки ІОР спрямований на визначення рівня сформованості професійних умінь, узагальнення результатів навчальної діяльності та оцінювання ефективності реалізованої підготовки. На цьому етапі відбувається зіставлення запланованих цілей із досягнутими результатами, що дає змогу виявити ступінь готовності здобувачів до розробки і впровадження ІОР у педагогічну практику. Метааналіз X.~Han та ін.~\cite{Han2025VRTeacherEd} підтверджує помірний позитивний ефект VR на результати підготовки вчителів (g=0,524) зі значними варіаціями залежно від рівня занурення й цілей навчання, що обґрунтовує багатокритеріальний підхід до оцінювання. Водночас SWOT-аналіз~\cite{Wang2024XRTeacherEd} ідентифікує <<відсутність конструктивного узгодження між цілями навчання, методами викладання й оцінюванням>> як ключову слабкість існуючих XR-програм підготовки вчителів --- рефлексивно-оцінний блок адресує цю проблему через явне зіставлення критеріїв із компонентами готовності.

Метою оцінно-результативного блоку є комплексне оцінювання результатів формування вмінь створювати імерсивні освітні ресурси з урахуванням технологічних, дидактичних та проєктних складників підготовки. У межах цього блоку визначаються критерії, показники та рівні сформованості відповідних умінь, здійснюється аналіз якості створених освітніх продуктів, а також оцінюється здатність майбутніх учителів інформатики обґрунтовувати педагогічну доцільність використання розроблених ресурсів.

До основних завдань оцінно-результативного етапу належить здійснення поточного й підсумкового контролю, організація самооцінювання та взаємооцінювання, аналіз динаміки професійного зростання студентів. Важливим компонентом цього етапу є рефлексія власної діяльності, що передбачає осмислення отриманого досвіду, виявлення труднощів і визначення шляхів подальшого вдосконалення професійних умінь.

Реалізація завдань цього блоку дозволяє не лише зафіксувати досягнуті результати, а й окреслити перспективи подальшого професійного вдосконалення майбутніх учителів інформатики у сфері розробки імерсивних освітніх ресурсів. Завершальним елементом структурно-функціональної моделі виступає результативний компонент, який відображає ефективність реалізації методики навчання розробки імерсивних освітніх ресурсів у процесі фахової підготовки майбутніх учителів інформатики та забезпечує систематичне осмислення студентами власного досвіду проєктування та створення ІОР, співвіднесення проміжних і підсумкових результатів із цілями навчання, а також оцінювання сформованості професійних умінь відповідно до визначених критеріїв і рівнів.

У результаті впровадження запропонованої методики формується готовність майбутніх учителів інформатики до розробки і педагогічно обґрунтованого використання імерсивних освітніх ресурсів, що виявляється у здатності методично доцільно добирати інструменти й технології, проєктувати структуру імерсивного освітнього продукту, інтегрувати його в освітній процес та здійснювати рефлексивне оцінювання власної педагогічної діяльності. Досягнення окресленого результату підтверджує методичну цілісність моделі, узгодженість її структурних блоків і результативність поетапної організації навчання розробки імерсивних освітніх ресурсів.

У межах дослідження означена готовність розглядається як інтегративна характеристика професійної підготовки, що відображає здатність до цілеспрямованого планування, організації та оцінювання процесу створення імерсивних освітніх ресурсів у навчальному середовищі та забезпечує методичну цілісність і результативність навчання розробки ІОР у процесі фахової підготовки майбутніх учителів інформатики.

Визначальну роль у структурі готовності відіграє \emph{\textbf{мотиваційно-ціннісний компонент},} який виконує смислоутворювальну й регулятивну функції та зумовлює спрямованість майбутнього вчителя інформатики на реалізацію методики навчання розробки імерсивних освітніх ресурсів. Його зміст виявляється у сформованості стійких професійних мотивів до педагогічно вмотивованого проєктування імерсивних освітніх продуктів, позитивному ставленні до використання цифрових технологій у навчальному процесі, готовності до безперервного професійного саморозвитку та усвідомленого вибору індивідуальної освітньої траєкторії. Саме цей компонент забезпечує методично обґрунтоване цілепокладання, визначення пріоритетів педагогічної діяльності та внутрішню настанову на досягнення результативності під час розв'язання професійних педагогічних і технологічних завдань. Емпірично засвідчено, що досвід VR-Making сприяє значному зростанню технологічної готовності та усвідомлення педагогічних переваг імерсивних технологій серед майбутніх учителів~\cite{Lee2022VRMaking}.

\emph{\textbf{Когнітивний компонент}} відображає систему професійно значущих знань, необхідних для реалізації методики навчання розробки імерсивних освітніх ресурсів, і формується у процесі опанування змісту фахових, методичних і психолого-педагогічних дисциплін. Він охоплює знання теоретичних засад імерсивного навчання, принципів педагогічного проєктування цифрових освітніх ресурсів, можливостей віртуальних і доповнених середовищ, а також дидактичних вимог до їх використання в освітньому процесі. Когнітивний компонент забезпечує усвідомлене прийняття методичних рішень і перехід від репродуктивного використання готових цифрових продуктів до самостійного педагогічного проєктування. Систематичний огляд~\cite{Mallek2024SynthesizingVR} підтверджує, що результати імерсивного навчання суттєво покращуються за умови узгодження навчального дизайну з теоріями конструктивізму та емпіричного навчання, а керування когнітивним навантаженням є критичною компетентністю розробника ІОР~\cite{Buchner2021CogLoad}.

\emph{\textbf{Діяльнісно-проєктувальний компонент}} характеризує рівень сформованості практичних умінь майбутніх учителів інформатики, пов'язаних із розробкою та впровадженням імерсивних освітніх ресурсів. Його зміст охоплює володіння апаратними й програмними засобами, уміння здійснювати відбір, структурування та педагогічну адаптацію навчального контенту, а також використання спеціалізованих програмних середовищ для створення віртуальних і доповнених освітніх ресурсів. У межах цього компонента формується здатність інтегрувати імерсивні ресурси в структуру уроку або навчального курсу з урахуванням освітніх цілей і методичних вимог. Практичне створення AR-ресурсів розвиває TPACK-компетентність~\cite{BeldaMedina2022TPACK}, а цикл емпіричного навчання забезпечує найвищу результативність формування проєктувальних умінь~\cite{Kee2023Experiential}.

\emph{\textbf{Технологічний компонент}} розкриває способи організації освітнього процесу з використанням цифрових та імерсивних технологій і відображає готовність до технологічного проєктування навчання. Він передбачає здатність планувати навчальну діяльність, визначати логіку подання навчального матеріалу, організовувати поетапне просування здобувачів освіти до досягнення запланованих результатів, а також добирати адекватні технологічні рішення відповідно до дидактичних завдань. Реалізація цього компонента забезпечує використання імерсивних освітніх ресурсів як інструменту методично виваженої педагогічної діяльності. Ключовим орієнтиром виступає <<інтерактивний дизайн зворотного зв'язку>>, визначений як найвпливовіший фактор ефективності VR-навчальних систем~\cite{Li2024ADDIE}, а також соціотехніко-педагогічна інтеграція, що враховує взаємозв'язок людського, технологічного та педагогічного компонентів~\cite{Schmidt2024EntangledLXD}.

\emph{\textbf{Інтерактивний компонент}} відображає здатність майбутнього вчителя інформатики організовувати педагогічну взаємодію в умовах цифрового й імерсивного освітнього середовища. Його зміст пов'язаний із розумінням інтерактивності як методичної характеристики освітнього процесу, що забезпечує активну участь здобувачів освіти, діалог, співпрацю та спільне розв'язання навчальних завдань. У межах використання імерсивних освітніх ресурсів цей компонент реалізується через створення умов для активної пізнавальної діяльності та педагогічно доцільної взаємодії з цифровим середовищем. Дослідження L.~Li та ін.~\cite{Li2024ClassMasterIVR} засвідчує, що повністю імерсивне VR-середовище забезпечує <<інноваційну, імерсивну та практичну>> підготовку, яка сприяє довготривалому збереженню компетентностей, а спільне проєктування IVR-ресурсів створює <<простір для практики й рефлексії>>~\cite{Paulsen2024CollaborativeIVR}.

\emph{\textbf{Рефлексивно-оцінний компонент}} завершує структуру готовності та забезпечує усвідомлення результатів власної педагогічної діяльності в контексті реалізації методики навчання розробки імерсивних освітніх ресурсів. Він виявляється у здатності оцінювати відповідність створених ресурсів дидактичним цілям, аналізувати ефективність їх використання в освітньому процесі, виявляти проблемні аспекти та здійснювати корекцію педагогічних і технологічних рішень, визначаючи напрями подальшого професійного розвитку. Метааналіз~\cite{Han2025VRTeacherEd} підтверджує, що VR сприяє формуванню рефлексивної практики у підготовці вчителів (g=0,524), а стадії рефлексивного спостереження (RO) й абстрактної концептуалізації (AC) циклу Д.~Колба є ключовими для перетворення імерсивного досвіду в стійкі професійні уміння~\cite{Kee2023Experiential}.

Компонентна структура готовності майбутніх учителів інформатики до реалізації методики навчання розробки імерсивних освітніх ресурсів узагальнено та подано на рис.~\ref{fig:2.4}.

\begin{figure}[!h]
\centering
\resizebox{\textwidth}{!}{%
\begin{tikzpicture}
\footnotesize

% Top header
\node[dia/rbox=17cm] (title) {\small\textbf{Готовність майбутніх вчителів інформатики до розробки імерсивних освітніх ресурсів}};

% Column headers (6 columns) — centers spaced at 2.85cm intervals
\node[dia/rbox=2.5cm, minimum height=1.4cm, below=1.5cm of title.south, xshift=-7.125cm, font=\footnotesize] (h1)
  {\textbf{Мотиваційно-ціннісний компонент}};
\node[dia/rbox=2.5cm, minimum height=1.4cm, below=1.5cm of title.south, xshift=-4.275cm, font=\footnotesize] (h2)
  {\textbf{Когнітивний компонент}};
\node[dia/rbox=2.5cm, minimum height=1.4cm, below=1.5cm of title.south, xshift=-1.425cm, font=\footnotesize] (h3)
  {\textbf{Діяльнісно-проєктувальний компонент}};
\node[dia/rbox=2.5cm, minimum height=1.4cm, below=1.5cm of title.south, xshift=1.425cm, font=\footnotesize] (h4)
  {\textbf{Технологічний компонент}};
\node[dia/rbox=2.5cm, minimum height=1.4cm, below=1.5cm of title.south, xshift=4.275cm, font=\footnotesize] (h5)
  {\textbf{Інтерактивний компонент}};
\node[dia/rbox=2.5cm, minimum height=1.4cm, below=1.5cm of title.south, xshift=7.125cm, font=\footnotesize] (h6)
  {\textbf{Рефлексивно-оцінний компонент}};

% Arrows from title to headers
\draw[dia/arrow] (title.south) -- ++(0,-0.4cm) -| (h1.north);
\draw[dia/arrow] (title.south) -- ++(0,-0.4cm) -| (h2.north);
\draw[dia/arrow] (title.south) -- ++(0,-0.4cm) -| (h3.north);
\draw[dia/arrow] (title.south) -- ++(0,-0.4cm) -| (h4.north);
\draw[dia/arrow] (title.south) -- ++(0,-0.4cm) -| (h5.north);
\draw[dia/arrow] (title.south) -- ++(0,-0.4cm) -| (h6.north);

% Detail boxes: all at same level below headers (no stagger)
\node[dia/lrbox=2.5cm, below=0.4cm of h1, font=\footnotesize] (d1) {%
  \textbullet~професійні мотиви;\\
  \textbullet~позитивне ставлення до технологій%
};

\node[dia/lrbox=2.5cm, below=0.4cm of h2, font=\footnotesize] (d2) {%
  \textbullet~теоретичні засади;\\
  \textbullet~дидактичні вимоги%
};

\node[dia/lrbox=2.5cm, below=0.4cm of h3, font=\footnotesize] (d3) {%
  \textbullet~володіння програмними засобами;\\
  \textbullet~інтеграція ресурсів%
};

\node[dia/lrbox=2.5cm, below=0.4cm of h4, font=\footnotesize] (d4) {%
  \textbullet~технологічне проєктування;\\
  \textbullet~педагогічна взаємодія%
};

\node[dia/lrbox=2.5cm, below=0.4cm of h5, font=\footnotesize] (d5) {%
  \textbullet~активна проєктувальна діяльність;\\
  \textbullet~співпраця%
};

\node[dia/lrbox=2.5cm, below=0.4cm of h6, font=\footnotesize] (d6) {%
  \textbullet~оцінка ефективності;\\
  \textbullet~корекція рішень%
};

% Arrows from headers to detail boxes
\draw[dia/arrow] (h1.south) -- (d1.north);
\draw[dia/arrow] (h2.south) -- (d2.north);
\draw[dia/arrow] (h3.south) -- (d3.north);
\draw[dia/arrow] (h4.south) -- (d4.north);
\draw[dia/arrow] (h5.south) -- (d5.north);
\draw[dia/arrow] (h6.south) -- (d6.north);

\end{tikzpicture}%
}

\caption{Структура готовності майбутніх учителів інформатики до реалізації методики навчання розробки імерсивних освітніх ресурсів}
\label{fig:2.4}
\end{figure}

Представлений рисунок відображає логіку взаємозв'язку мотиваційно-ціннісного, когнітивного, діяльнісно-проєктувального, технологічного, інтерактивного та рефлексивно-оцінного компонентів, які в сукупності забезпечують методичну цілісність процесу підготовки майбутніх учителів інформатики.

Зв'язок між розробленою моделлю та практичною діяльністю автора забезпечується конкретною реалізацією кожного блоку: мотиваційно-цільовий блок втілено через WebAR-квести для профорієнтації~\cite{Shepiliev2021Career}, що демонструють практичну цінність імерсивних технологій; змістовий блок реалізовано у курсі <<Синтетична реальність в освіті>> (18 модулів) на основі навчального посібника з WebAR; організаційно-технологічний блок --- через прогресивну методику від базового відстеження зображень до ML-інтегрованого WebAR~\cite{Semerikov2024WebARML_CoSinE, Semerikov2024WebAR_EduDim}; рефлексивно-оцінний блок --- через багатокритеріальне оцінювання, узгоджене з п'ятьма характеристиками ІОР (підрозділ~\ref{subsec:1.2}). Ця відповідність забезпечує послідовний перехід від теоретичної моделі до практичного впровадження та експериментальної перевірки (розділ~3).

Розроблена структурно-функціональна модель підготовки майбутніх учителів інформатики до розробки імерсивних освітніх ресурсів забезпечує цілісне осмислення методики формування відповідних професійних умінь у процесі фахової підготовки. Її реалізація уможливлює системне поєднання цільових орієнтирів, змістових компонентів, організаційно-технологічних механізмів і результативних показників, а також дозволяє окреслити методичні можливості й обмеження організації освітнього процесу. Практичне впровадження моделі створює передумови для цілеспрямованого розвитку здатності майбутніх учителів інформатики до педагогічно вмотивованого та технологічно обґрунтованого проєктування імерсивних освітніх ресурсів у професійній діяльності.
